\chapter{LITERATURE REVIEW: PARTIAL DISCHARGE DETECTION USING TRADITIONAL MACHINE LEARNING}

\section{Machine Learning in Partial Discharge Detection}

The application of machine learning (ML) techniques to partial discharge (PD) analysis has become increasingly prevalent, driven by the need to overcome the limitations inherent in conventional diagnostic methods \parencite{Lu_2020_2, Masud_2016}. Historically, PD assessment relied heavily on the manual interpretation of measurement data, such as Phase-Resolved Partial Discharge (PRPD) patterns, by human experts. This process is often subjective, time-consuming, and can be unreliable, particularly when dealing with the large volumes of data generated by modern online monitoring systems \parencite{Alsumaidaee_2022, Seri_2021}. Traditional approaches also face significant challenges in noisy operational environments, where distinguishing PD signals from external interference and background noise is a primary obstacle \parencite{Ghosh_2020, Park_2023_1}. Moreover, identifying and separating signals from multiple concurrent PD sources \parencite{Baug_2017_1, Firuzi_2017_1} and reliably quantifying the severity or danger of an insulation defect remain complex tasks for conventional analysis \parencite{Ma_2021_3, Orellana_2023}. Machine learning offers a pathway to automated, objective, and more accurate PD analysis, forming the core of advanced condition monitoring and diagnostic systems \parencite{Lu_2020_2}. A primary application of ML in this field is the classification of PD source type such as differentiating between corona, surface discharges, and internal voids which is essential for assessing the risk posed to insulation integrity \parencite{Anjum_2017_1, Borghei_2022}. Another key task is anomaly detection, where ML models are trained to identify PD events as deviations from normal operational signals or ambient noise, a capability crucial for robust online monitoring \parencite{Florkowski_2021_1, Thi_2020}. Furthermore, ML algorithms are increasingly used for the localization of PD sources, estimating their physical position within high-voltage (HV) equipment, which facilitates targeted maintenance and repair \parencite{Chan_2023, Liu_2020_4}. The overarching motivation for employing ML is to develop intelligent systems that can learn complex patterns directly from sensor data, adapt to changing operational conditions, and provide consistent and reliable diagnostic assessments \parencite{Sikorski_2025}.

\section{Overview of Machine Learning Algorithms for Partial Discharge Detection}

A variety of traditional machine learning algorithms have been employed for Partial Discharge (PD) detection and analysis. These algorithms are primarily used for tasks such as the classification of PD types, localization of PD sources, and noise rejection. This section provides an overview of some commonly used traditional ML algorithms found in the reviewed literature.

\subsection{Support Vector Machines (SVM)}

Support Vector Machines (SVMs) are supervised learning models that analyze data for classification and regression. In classification tasks, SVMs construct an optimal hyperplane in a high-dimensional space to serve as a decision boundary with the largest possible margin to the nearest data points of any class. SVMs have been widely used for classifying PD sources and defect types. For instance, SVM has been effectively utilized to recognize PDs in switchgear from ultra-high frequency (UHF) signal data \parencite{Jin_2018_1} and to classify four different PD types using an optimized model \parencite{Dai_2017_1}. They are also used for localizing PD sources in GIS based on image edge detection features \parencite{Li_2019_7}. Advanced applications include using a Brain Storm Optimization (BSO) enhanced SVM for pattern recognition of ultrasonic PD signals in transformers \parencite{Zhou_2022_4}. The selection of SVM is often motivated by its strong theoretical foundation, effectiveness in high-dimensional spaces, and robust performance even with limited training data \parencite{Bin_2020_1, Zhou_2017_3}. Its ability to handle the complex characteristics of PD data and its strong performance in multi-class classification make it a preferred choice for pattern recognition and distinguishing between different PD types \parencite{Jin_2018_1}, Sertta? et al., 2020. Furthermore, One-Class SVM (OCSVM) has been chosen for its ability to model normal, PD-free conditions without requiring a fully labeled training set, which is advantageous for online monitoring systems operating in low signal-to-noise ratio (SNR) environments Parrado-Hern?ndez et al., 2018. For defect classification in ceramic insulators, SVM outperformed k-NN, Naive Bayes, and Decision Tree with 90.6\% accuracy \parencite{Ma_2021_3}. For aging state recognition in polyethylene-insulated cables, SVM achieved up to 100\% accuracy with a small subset of features \parencite{Morette_2020_1}.

\subsection{Random Forest (RF)}

Random Forest (RF) is an ensemble learning method that constructs a multitude of decision trees during training. It outputs the class selected by the most trees for classification tasks, making it highly regarded for its robustness against overfitting and its ability to manage high-dimensional data. RF has been successfully applied in various PD detection scenarios, including identifying the location of PD sources \parencite{Bag_2019_1} and classifying PD types \parencite{Guo_2021, Prabhu_2016_1}. Specific applications include locating PD sources using features from optical sensor data \parencite{Bag_2019_1}, identifying faults based on the concentration of decomposition gases \parencite{Tan_2022}, and classifying defects in on-board cable terminals for high-speed railways \parencite{Guo_2021}. The choice of RF is often driven by its high accuracy, its ability to process large datasets with numerous input variables, and its built-in feature ranking capabilities \parencite{Bag_2019_1}. In terms of performance, RF has consistently demonstrated high accuracy. A PD source location accuracy of 95.6\% was achieved using RF on optical sensor data \parencite{Bag_2019_1}. \subsection{Decision Trees (DT)}

Decision Trees are a non-parametric supervised learning method used for classification and regression. They build a model that predicts a target variable's value by learning simple, interpretable decision rules inferred from data features. Decision trees have been applied for PD fault type recognition \parencite{Lee_2020_1} and for detecting insulation faults in transformers, where they have shown high accuracy for specific fault classes \parencite{Maseko_2025}. These models are often chosen for their interpretability and their capacity to handle both numerical and categorical data. They have also served as a baseline model for comparison against more complex algorithms like Back Propagation Neural Networks (BPNN) \parencite{Fei_2024_1}. In terms of performance, a Fuzzy Logic Clustering Decision Tree (FLCDT) was shown to perform significantly better than standard decision tree algorithms like CART and See5 in terms of classification precision for transformer diagnostics \parencite{Lee_2020_1}. Furthermore, an accuracy of 100\% was achieved for specific insulation faults, such as those caused by thermal aging, using \parencite{Maseko_2025_1}.

\subsection{k-Nearest Neighbors (k-NN)}

k-Nearest Neighbors (k-NN) is a simple, instance-based learning algorithm. For classification, an object is assigned to the class most common among its k k nearest neighbors, as determined by a majority vote. k-NN has been applied to classify PD types in GIS \parencite{Bin_2020_1, Fei_2024_1} and for online stator winding fault diagnosis in induction motors \parencite{Sarkar_2021}. The algorithm is frequently chosen for its simplicity, effectiveness, and ease of implementation, especially in scenarios with irregular decision boundaries \parencite{Bin_2020_1}. It has also been noted for enabling effective online fault detection with high prediction accuracy for motor faults \parencite{Sarkar_2021}. In terms of performance, k-NN achieved an accuracy of 88.5\% in classifying four types of PD in GIS \parencite{Bin_2020_1}. In motor fault diagnosis, it demonstrated acceptable prediction accuracy, making it viable for online monitoring systems \parencite{Sarkar_2021}. While serving as a baseline model, k-NN's results have been compared against ANNs that achieved classification accuracies as high as 96.03\% \parencite{Haiba_2024}.

\subsection{Artificial Neural Networks (ANN) (Traditional Architectures)}

Artificial Neural Networks (ANNs), in their traditional forms like multilayer perceptrons (MLPs) trained with back-propagation, are computational models inspired by the brain's structure. These models consist of interconnected nodes arranged in layers. Other traditional ANN architectures used in PD detection include Self-Organizing Maps (SOM), which are used for clustering and visualization, and Learning Vector Quantization (LVQ) networks. ANNs have been widely applied for PD pattern recognition and classification of defect types. For example, ANNs have been used to classify discharges from different defects in ceramic insulators \parencite{Anjum_2017_1} and to diagnose faults in motor stator windings \parencite{Chen_2018_3}. A Back-Propagation ANN (BP-ANN) has been used for PD identification in GIS based on acoustic emission signals \parencite{Hou_2021_2}, and a Learning Vector Quantization (LVQ) network has been used for localizing PD in noisy environments by analyzing winding sections \parencite{Mohammadirad_2025}. ANNs are selected for their ability to model complex, non-linear relationships between input features and output classes, as well as their capacity for learning directly from data \parencite{Chen_2018_3, Hou_2021_2}. They have demonstrated high recognition rates for various PD faults and the capability to determine degradation stages Mas'ud et al., 2016. In terms of performance, recognition rates exceeding 95\% were reported for each defect class using ANNs \parencite{Anjum_2017_1}. An accuracy of 90\% was achieved in diagnosing stator faults \parencite{Chen_2018_3}, and 96.03\% accuracy was obtained for insulator defect classification using discrete wavelet transform features \parencite{Haiba_2024}. A study classifying PD types reported a 97\% recognition rate \parencite{Zahed_2017_1}, and a BP-ANN model for GIS defects achieved 93.7\% accuracy \parencite{Hou_2021_2}.

\subsection{Fuzzy Logic Based Classifiers}

Fuzzy logic systems operate on "degrees of truth" rather than binary true/false logic, making them well-suited for managing the uncertainty and imprecision inherent in real-world data. Common fuzzy logic-based classifiers include Fuzzy Decision Trees, Fuzzy C-Means (FCM) clustering, and Adaptive Neuro-Fuzzy Inference Systems (AN- FIS). Fuzzy logic has been applied to diagnose PD in cast-resin transformers \parencite{Lee_2020_1} and for rejecting interference \parencite{Wang_2019_3}. It is also used for identifying PD signals through clustering \parencite{Chen_2017_2}. An Adaptive Neuro- Fuzzy Inference System (ANFIS), which integrates neural networks with fuzzy logic, has been used for evaluating PD intensity by leveraging its ability to continuously optimize model parameters and reduce errors compared to traditional methods \parencite{Wang_2019_3}. The adaptability and robustness of ANFIS make it a strong candidate for PD analysis. Fuzzy logic methods are particularly valuable because they can be designed to provide interpretable, rule-based outputs, which aids in understanding the diagnostic decision-making process. Regarding performance, the Fuzzy Logic Clustering Decision Tree (FLCDT) was found to perform significantly better than traditional decision trees in terms of classification precision \parencite{Lee_2020_1}. An adaptive Fuzzy C-Means (FCM) method reported high accuracy in identifying PD signals \parencite{Chen_2017_2}. Furthermore, an improved ANFIS model showed a 2\% reduction in model error compared to a traditional ANFIS when evaluating PD intensity \parencite{Wang_2019_3}.

\subsection{Clustering Algorithms (K-Means, GMM, DBSCAN, etc.)}

Clustering algorithms are unsupervised learning techniques that group objects such that those in the same cluster are more similar to each other than to those in different clusters. Common methods include K-Means, Gaussian Mixture Model (GMM), and density-based approaches like DBSCAN. Self-Organizing Maps (SOM) can also be viewed as a clustering technique. These algorithms have been applied to separate signals from multiple PD sources \parencite{Firuzi_2017_1, Mishra_2019_1} and perform data cleaning by separating signal from noise \parencite{Soh_2022}. They can also identify distinct PD patterns for localization \parencite{Zhou_2021_4}. Methods like K-Means, GMM, and Mean-shift have been combined to classify PD signals effectively \parencite{Carvalho_2024}. Self-Organizing Maps, when paired with statistical feature transformation, have been employed for PD classification \parencite{Bohari_2021}. A notable advancement is the development of PDAutoClust, a parameter-free clustering technique designed to reduce reliance on user-defined parameters, making the process more robust and accessible \parencite{Rahman_2024_1}. Clustering algorithms are valuable because they can discover underlying structures in PD data without requiring pre-labeled datasets. This is crucial for unsupervised monitoring where fault types may not be known in advance. GMM clustering helps discover distinct groups in a feature space for separating mixed PD signals \parencite{Firuzi_2017_1}, while density-based methods can effectively separate PD pulses from various types of noise and interference \parencite{Zhou_2021_4}. Performance-wise, a combination of K-Means, GMM, and Mean-shift clustering followed by an SVM classifier achieved 93\% accuracy \parencite{Carvalho_2024}. The PDAutoClust method performed favorably against other density-based algorithms on PD datasets \parencite{Rahman_2024_1}. A SOM-based classifier demonstrated classification times under five seconds with nearzero error rates \parencite{Bohari_2021}, and an advanced density peak clustering algorithm reported an average localization error of just 0.21 m for two simultaneous PD sources \parencite{Zhou_2021_4}.

\section{Feature Engineering in Partial Discharge Analysis}

\subsection{Time-Domain Features}

Time-domain features are derived directly from the individual PD pulse waveforms to quantify their fundamental physical properties. These features characterize the shape, magnitude, and temporal occurrence of the pulse, which are intrinsically linked to the underlying discharge physics and provide crucial information for classification \parencite{Lu_2020_2}. The analysis of these characteristics enables the differentiation of various PD phenomena based on how the discharge manifests as a transient electrical signal \parencite{Boczar_2017}. The peak amplitude, signifying the maximum magnitude of the pulse, serves as a primary indicator of the discharge intensity and is directly related to the energy released during the event \parencite{Zhang_2024_28}. Different PD sources, such as corona, surface discharges, or internal voids, often generate pulses with distinct amplitude ranges. This makes amplitude a valuable feature for assessing a defect's severity and for distinguishing between source types, as the physical mechanism and geometry of the defect govern the amount of charge mobilized in the discharge \parencite{Boczar_2017}. The pulse rise time, defined as the interval for the signal to increase from 10\% to 90\% of its peak value, reflects the speed of the ionization process \parencite{Lu_2020_2}. PD pulses are characterized by very rapid rise times, often in the nanosecond or subnanosecond range \parencite{Darwish_2019}. This parameter is highly effective for discriminating between PD sources because phenomena such as streamer-like corona discharges exhibit different development speeds compared to Townsend-like discharges in internal voids \parencite{Xavier_2021}. Consequently, rise time is a critical feature, though it can be affected by signal degradation and dispersion during propagation from the source to the sensor \parencite{Lu_2020_2}. Pulse duration, frequently measured as the full width at half maximum (FWHM), quantifies the time for which the discharge event persists and is related to the recombination time of charge carriers following the initial ionization Mart?nez-\parencite{MartinezTarifa_2022}. This temporal characteristic is another key descriptor of the pulse shape that aids in distinguishing between defect types. Specific discharge mechanisms, influenced by factors like gas pressure and electrode geometry, produce characteristically longer or shorter pulses, making duration a useful parameter for classification \parencite{Boczar_2017, Uwiringiyimana_2022}. Finally, the pulse count per cycle tallies the total number of PD events detected within one period of the power frequency, directly quantifying the repetition rate and activity level of the PD source \parencite{Gillis_2023}. A high pulse count can indicate a more active or severe insulation defect, and monitoring this metric is essential for tracking the evolution and stability of a fault over time \parencite{Ma_2021_3}. The number of discharges per cycle, along with their phase distribution, is fundamental to forming Phase-Resolved Partial Discharge (PRPD) patterns, which are widely used for identifying specific defect types \parencite{Boczar_2017, Kasinathan_2017}.

\subsection{Statistical Features}

Statistical features provide a quantitative description of the shape and distribution of discharge activity within the Phase-Resolved Partial Discharge (PRPD) pattern. These descriptors are derived from the distribution of pulse magnitudes or counts across the phase angles of the AC voltage cycle. For this work, the higher-order statistical moments of skewness and kurtosis were computed, as they are particularly effective for capturing the nuanced characteristics of PRPD patterns that differentiate various discharge phenomena \parencite{Boczar_2017}. Skewness measures the degree of asymmetry in the distribution of PD pulses relative to the center of the phase axis. This feature is valuable because the physical mechanisms governing PD can differ significantly between the positive and negative half-cycles of the applied voltage, leading to characteristically asymmetric PRPD patterns \parencite{Boczar_2017}. For instance, corona from a sharp point electrode typically produces a much higher discharge activity during one polarity of the AC cycle, resulting in a highly skewed distribution \parencite{Yao_2016}. By quantifying this asymmetry, skewness provides a powerful metric for distinguishing between sources like corona, surface, and internal void discharges, each of which exhibits unique symmetry properties in its PRPD profile \parencite{Bohari_2021}. Kurtosis quantifies the "peakedness" or "tailedness" of the distribution, indicating whether the PD pulses are concentrated within narrow phase windows or are more broadly dispersed across the cycle. A high kurtosis value signifies a distribution with sharp, narrow peaks and heavy tails, which is characteristic of discharge types that occur only at specific, repeatable moments near the voltage peaks \parencite{Boczar_2017}. Conversely, a lower kurtosis suggests a flatter, more uniform distribution of pulses. Since different insulation defects generate discharges with distinct phase concentrations, kurtosis serves as a key feature for identifying the nature of the PD source by characterizing the shape of its statistical distribution \parencite{Kunicki_2018}. Together, these statistical moments transform the qualitative visual information of a PRPD pattern into a compact and discriminative numerical feature vector suitable for automated classification \parencite{Bohari_2021}.

\subsection{Frequency-Domain Features}

Frequency-domain analysis is employed to investigate the spectral composition of partial discharge signals, which provides critical diagnostic information not apparent in the time domain. By applying a Fast Fourier Transform (FFT) to each captured waveform, the signal is converted from its temporal representation to its constituent frequencies \parencite{Li_2018_2}. The significance of this analysis lies in the fact that PD pulses are inherently fast, transient events with rapid rise times, meaning they are broadband phenomena containing significant energy across a wide spectrum that can extend into the High Frequency (HF), Very High Frequency (VHF), and Ultra-High Frequency (UHF) ranges \parencite{Darwish_2019}. The frequency signature of a measured PD signal is shaped by two key factors: the physical nature of the discharge source and the propagation path the electromagnetic wave travels to the sensor \parencite{Azirani_2021}. The equipment itself, such as a power transformer or a Gas-Insulated Switchgear (GIS) enclosure, acts as a complex waveguide, causing frequency-dependent attenuation and reflection that alters the signal's spectral characteristics \parencite{Promwong_2020}. Consequently, the energy distribution across different frequency bands becomes a unique fingerprint for a specific defect type within a given apparatus. Different types of insulation defects are known to manifest distinctively in the frequency domain. The UHF band (300 MHz--3 GHz) is particularly valuable for PD detection in transformers and GIS, as signals in this range can propagate efficiently within the metallic enclosures while being less susceptible to external corona and broadcast interference \parencite{Salah_2022}. For instance, the spectral content of UHF signals can be analyzed to diagnose and distinguish between specific fault types in GIS, such as those caused by free-moving particles or fixed protrusions on conductors \parencite{Ren_2021_1}. Selecting an optimal frequency band for analysis is therefore crucial for maximizing the signal-to-noise ratio and detection sensitivity \parencite{Azirani_2021}. For this study, the spectral power was computed within the distinct HF, VHF, and UHF bands to create a feature that captures this discriminative energy distribution, providing a robust basis for classification \parencite{Maneerot_2019}.

\subsection{Time-Frequency Features}

Partial discharge signals are fundamentally transient and non-stationary phenomena, meaning their frequency content evolves rapidly over their short duration. This characteristic makes time-frequency analysis methods, which can capture both temporal and spectral information simultaneously, particularly well-suited for their characterization \parencite{Li_2018_2}. Among these methods, the Discrete Wavelet Transform (DWT) is exceptionally effective and widely employed in PD analysis for its dual capability of signal denoising and feature extraction \parencite{Lu_2020_2}. The DWT decomposes a signal into various frequency sub-bands at different resolutions, allowing for the effective isolation of PD-related components from background noise and interference \parencite{Sahnoune_2024}. The efficacy of the DWT is highly dependent on the choice of the mother wavelet, as its shape should ideally correlate with the morphology of the PD pulses being analyzed to achieve optimal decomposition \parencite{Sahnoune_2024}. For this study, the Daubechies 4 (db4) wavelet was selected. The Daubechies family of wavelets is frequently chosen for its properties of compact support and orthogonality, and the db4 wavelet, in particular, has been demonstrated in the literature to be effective for the denoising and analysis of PD signals, thereby representing an established and validated choice \parencite{Wang_2020_10}. The signal was decomposed into four levels, which separates the signal into a series of detail coefficients representing successively narrower high-frequency bands, and a final approximation coefficient representing the low-frequency content. For feature extraction, the energy of the detail coefficients at each of the four decomposition levels was calculated. The resulting four-dimensional vector of wavelet energies serves as a powerful feature set. Since different types of PD sources (e.g., surface, corona, void) generate pulses with distinct spectral signatures, their energy is distributed differently across these frequency bands \parencite{Gao_2023_1}. \section{Dataset Diversity and Model Generalizability}

The generalizability of machine learning (ML) models for Partial Discharge (PD) detection is fundamentally challenged by the immense diversity of PD data \parencite{Lu_2020_2}. Signal characteristics are highly sensitive to the defect type \parencite{Lee_2020_1, Zang_2022_1, Bin_2020_1, Meitei_2020}, and operational environment. Consequently, a model trained on a specific dataset, particularly laboratory-generated data, often fails to generalize to different realworld scenarios \parencite{Raymond_2022, Sun_2019_1}. To mitigate this, researchers employ strategies like domain adaptation \parencite{Niu_2025, Tai_2024, Tuyet_Doan_2024}, and data augmentation to improve model robustness \parencite{Rauscher_2024, Raymond_2022}. Despite these efforts, a critical gap persists: many studies still rely on limited, homogeneous datasets from controlled lab settings \parencite{Zhang_2016_5}. The scarcity of real-world data and the absence of standardized public benchmarks severely hinder the development of truly robust and widely applicable ML solutions \parencite{Kim_2020_1, Lu_2020_2}.

\section{Comparative Studies and Benchmarking in ML for Partial Discharge Detection}

Comparative studies and benchmarking are essential for identifying the most effective algorithms and methodologies in ML-based PD detection. Several studies have undertaken such comparisons, revealing performance differences among models. For example, \parencite{Bin_2020_1} found that a PSO-ELM outperformed SVM and k-NN for GIS defect classification. \parencite{Lee_2020_1} showed their Fuzzy Logic Clustering Decision Tree surpassed standard decision trees for transformer diagnostics. Likewise, \parencite{Ma_2021_3} demonstrated that SVM was superior to k-NN, Naive Bayes, and Decision Trees for classifying PD severity. However, existing comparative analyses often suffer from inconsistencies that make direct, field-wide comparisons challenging \parencite{Lu_2020_2}. A primary issue is the use of different, often private, datasets, which limits the generalizability of findings \parencite{Kim_2021_1}. Since PD signal characteristics vary significantly with the defect type, insulation, and equipment, results from one dataset are not directly applicable to another \parencite{Ranjan_2024}. A key limitation is that comparisons often focus on algorithmic performance without deeply exploring how feature engineering strategies interact with the choice of model \parencite{Govindarajan_2021_2}. The optimal ML algorithm for a given task might change depending on the features extracted from the raw signal. There is a notable lack of standardized benchmarking protocols that test across various feature-algorithm pairings and generalization scenarios \parencite{Lu_2020_2}. This need for standardization was highlighted in the development of a qualification procedure for PD analysers \parencite{Vera_2023}. Addressing these issues is essential for developing truly robust and reliable ML-based PD detection systems.

\section{Research Gaps}

Despite the promise of machine learning (ML) for Partial Discharge (PD) detection, a review of the literature reveals critical gaps that hinder the development of robust and generalizable diagnostic tools. A primary gap is the limited exploration of featurealgorithm interactions; most studies evaluate ML models using a single, predefined set of features \parencite{Lu_2020_2}, neglecting to systematically investigate how different feature combinations impact the performance of various algorithms \parencite{Govindarajan_2021_2, Nimmo_2017_1}. Furthermore, many models are trained and tested on limited, homogeneous datasets from a single equipment type or condition \parencite{Kim_2021_1, Woon_2016_1}. This practice raises significant questions about their ability to generalize to unseen PD sources and diverse operational environments, a key barrier to practical deployment \parencite{Lu_2020_2}. This research addresses these gaps by systematically comparing traditional ML models across diverse PD datasets. We will evaluate how various feature combinations affect each algorithm's performance to identify optimal pairings, a need highlighted by the lack of such analysis in current literature \parencite{Govindarajan_2021_2}. Furthermore, we will explicitly test model generalizability using a combined dataset reflecting real-world conditions from various equipment like transformers, GIS, and cables \parencite{Bin_2020_1, Meitei_2020}, aiming to develop more robust and reliable diagnostic tools.
